% !TeX spellcheck = ru_RU
% !TEX root = vkr.tex

\section{Обзор}
\label{sec:relatedworks}

В данном разделе нужно описать всё, что необходимо для понимания Вашей работы и что придумали не Вы.
В дальнейших разделах нельзя прерывать повествование, например, для рассказа о деталях используемой технологии или архитектуре старой системы, потому что читателю будет трудно отличить Ваш вклад от не Вашего.

Любой обзор пишется с какой-то целью (обосновать актуальность, найти и описать интересные решения, сравнить и выбрать технологии) и по какой-то методике поиска материала (например, поиск N релевантных статей на таких-то сервисах).
Не будет лишним это всё явно описать.

\subsection{Обзор существующих решений}

В настоящее время существует ряд программных средств и библиотек,
предназначенных для моделирования мультиагентных систем и трекинга объектов. Однако большинство из них либо
ориентированы на централизованные методы обработки, либо не учитывают
специфику рассматриваемой задачи — анализ мультиагентных алгоритмов
в условиях возмущений и помех, не обладающих стандартными статистическими
предположениями.

\textbf{MATLAB Sensor Fusion and Tracking Toolbox} предоставляет широкий
набор инструментов для моделирования сенсоров, трекинга объектов и оценки
качества алгоритмов. К его преимуществам относятся развитая экосистема,
наличие готовых реализаций алгоритмов и удобные средства визуализации.
Однако данный инструмент в первую очередь ориентирован на централизованные
подходы к обработке данных. Кроме того, система является закрытой и требует
дорогостоящей лицензии, а интеграция пользовательских мультиагентных
алгоритмов затруднена.

\textbf{Repast (Agent-Based Modeling Toolkit)} представляет собой универсальную
платформу для моделирования мультиагентных систем. К её преимуществам можно
отнести гибкость и возможность описания широкого класса агентных моделей.
Тем не менее, данная платформа в основном применяется в социально-экономических
и биологических задачах и не содержит встроенных средств для моделирования
сенсоров и алгоритмов трекинга. Также стоит отметить, что система в настоящее
время активно не развивается и ограниченно поддерживается.

\textbf{SIMLAN и симуляторы на базе Gazebo} широко используются в задачах
робототехники для моделирования поведения роботов и тестирования алгоритмов
навигации. Их основным преимуществом является высокая степень реалистичности
физической симуляции и развитые средства интеграции с робототехническими
платформами. Однако такие системы ориентированы на детальное моделирование
физической среды и являются избыточными и сложными в настройке для задач
абстрактного анализа алгоритмов трекинга. Кроме того, они не предоставляют
удобных средств для массового сравнительного анализа алгоритмов.

\vspace{0.5em}

Для наглядного сравнения рассмотренных решений приведём таблицу
(табл.~\ref{tab:comparison_existing_systems}).

\begin{table}[h]
    \centering
    \caption{Сравнение существующих решений}
    \label{tab:comparison_existing_systems}
    
    \resizebox{\textwidth}{!}{
    \begin{tabular}{|l|c|c|c|}
    \hline
    \textbf{Характеристика} & \textbf{MATLAB Toolbox} & \textbf{Repast} & \textbf{Gazebo / SIMLAN} \\
    \hline
    Поддержка мультиагентных алгоритмов & $\pm$ & + & $\pm$ \\
    \hline
    Встроенные модели сенсоров & + & - & + \\
    \hline
    Инструменты трекинга & + & - & $\pm$ \\
    \hline
    Простота интеграции своих алгоритмов & - & + & - \\
    \hline
    Пакетные эксперименты и сравнение & $\pm$ & - & - \\
    \hline
    Визуализация результатов & + & $\pm$ & + \\
    \hline
    Ориентация на распределённые алгоритмы & - & $\pm$ & $\pm$ \\
    \hline
    Нестандартные статистические условия & - & - & - \\
    \hline
    \end{tabular}
    }
    \end{table}

\vspace{0.5em}

\textbf{Вывод.}
Проведённый анализ показал, что существующие решения либо ориентированы
на централизованные методы обработки данных, либо не предоставляют
необходимых инструментов для моделирования сенсорных систем и трекинга,
либо избыточны для рассматриваемой задачи. Кроме того, ни одна из систем
не поддерживает в полной мере проведение экспериментов в условиях
нестандартных статистических предположений.

Таким образом, существует необходимость в разработке специализированной
системы, позволяющей интегрировать пользовательские мультиагентные алгоритмы,
гибко задавать параметры симуляции и проводить их сравнительный анализ
в условиях возмущений и помех произвольной природы.

